\documentclass[11pt,a4paper]{article}
\usepackage[T1]{fontenc}
\usepackage{lmodern,amsmath,amssymb,amsthm,mathtools,microtype}
\usepackage[margin=27mm]{geometry}
\usepackage[colorlinks=true,linkcolor=blue,citecolor=blue,urlcolor=blue]{hyperref}
\setlength{\emergencystretch}{2em}
\newtheorem{theorem}{Theorem}[section]
\newtheorem{lemma}[theorem]{Lemma}
\newtheorem{corollary}[theorem]{Corollary}
\newcommand{\Q}{\mathbb Q}\newcommand{\Z}{\mathbb Z}
\DeclareMathOperator{\Tr}{Tr}
\title{Unbounded auxiliary poles with fixed elliptic detector roots}
\author{}\date{17 November 2026}
\begin{document}\maketitle
\begin{abstract}
In the complete normalized twenty-five-coordinate source, the heights
of the formal comparison ratios are not bounded by the fixed root and
contact data. We construct complete rational source tuples with these
data fixed and with both rational components of one comparison ratio
having arbitrarily deep poles at one fixed prime. The construction uses
the source's retained prime-pattern and intrinsic pole-realization
premises. It concerns the existing detector equations; an additional
Pell-cover equation may impose further restrictions.
\end{abstract}

\section{Source and scope}

Use the preferred scalar-rescaled intrinsic source of
\cite{Scalar}, with the order-two comparisons of \cite{Precision}.
The same argument applies to its order-three comparisons.
Its curve, point roles and comparison function are
\[
 K=\Q(\sqrt2),\qquad E:y^2=x^3-(11-6\sqrt2)x,
\]
\[
 A,S,T_1,Z,H,U_A,U_S,U_T,U_Z,
 \qquad \ell_8(Q)=-\frac{x(Q)}{y(Q)}
       \left(1-\frac{2(11-6\sqrt2)}{5x(Q)^2}\right).
\]
Write the four root points as \(Q_i=[n_i]P\), with
\(Q_S=S\). They are fixed multiples of the retained non-torsion
point \(P\). Once these roots and the rational slope are fixed, the
two physical weights
\[
 w_j=\frac{\rho^*R_j}{(1+\rho^*R_j)^2}
\]
are fixed rational numbers. The remaining detector inputs are
\[
 V_j=\Tr(W_j^{-1}),\qquad b_j=4w_j/V_j^2.
\]
The intrinsic pole block is
\[
 Y_j^2=b_j(b_j+v_j^2)^2s_j^4
          -b_j(b_j-v_j^2)^2(1-b_j)^2.
\]
Its prepared completeness is used in precisely the scope in
\cite{Scalar,Intrinsic}. All original point, difference, scalar and
detector guards are retained.

The source has no equation requiring \(U_i=[h]Q_i\) when
\(H=[h]P\). This equality is a choice in its completeness proof.
After one prepared tuple has been constructed, the rational source
therefore permits its five auxiliary points to vary independently.

\begin{lemma}\label{lem:even}
Fix a nonzero rational weight \(w\). At every rational prime where
\(v_p(w)\) is even, the valuation of \(4w/V^2\) is even for
every nonzero rational \(V\). In particular, the square-denominator
condition imposes no bound on the depths of poles of the detector
functions at such a prime.
\end{lemma}
\begin{proof}
Its valuation is \(2v_p(2)+v_p(w)-2v_p(V)\), which is even.
\end{proof}

\section{A fixed unused prime}

We recall two explicit consequences of the retained simultaneous
prime-pattern premise. One may append finitely many fixed rational
contact coefficients before the pattern is chosen, prescribe avoidance
of any fixed finite set of primes, and make the positive slope
arbitrarily large. The source also supplies compatible primary and
opposite reduction-order conditions on the reference index.

\begin{lemma}\label{lem:dummy}
In the ordered completeness construction, one can select an additional
split prime \(p_*\), used by neither detector, with the following
properties:
\begin{enumerate}
\item Both physical weights are units at \(p_*\).
\item The values \(\ell_8(P)\) and \(\ell_8(S)\), the nonzero
root index \(n_S\), and all fixed good-reduction data are units at
both places above \(p_*\).
\item Every selected rational prime is coprime to all the prescribed
primary and opposite reduction orders. The prime \(p_*\) also avoids
the modulus for the previously prepared finite and real neighborhoods.
\item The initial reference index \(h_0\) can satisfy all the usual
completeness conditions and \(v_{p_*}(h_0)=1\).
\end{enumerate}
\end{lemma}
\begin{proof}
Fix the integer input, scalar, roots, scales and preparation
neighborhoods in the original order. Choose a positive rational
coefficient \(c_*\) distinct from every actual contact and pole
coefficient, such that \(R_j(c_*)\) and
\(1+\rho^*R_j(c_*)\) are nonzero for \(j=0,1\).
There are only finitely many forbidden values, and sufficiently large
positive values have these properties. Enlarge the fixed avoidance
list by the supports of these rational values, of \(\rho^*\), of
\(n_S\), and of the nonzero fixed point-function values in item two.
All these data precede prime selection.

Append \(c_*\) to the prime pattern. In rational-contact notation,
the additional prime is \(p_*=u-c_*v\), where \(\lambda=u/v\),
with the prescribed fixed denominator clearing understood. The chosen
prime avoids that denominator. Thus \(\lambda\equiv c_*\pmod{p_*}\),
which proves item one. The avoidance conditions prove item two.

For completeness, the elementary separation used by the source applies
at both conjugate places. Take the slope sufficiently large that every
selected prime lies in \((3u/4,u)\), with \(u\ge25\).
The full elliptic reduction order at either place above any selected
prime is even and less than
\(u+1+2\sqrt u<3u/2\). A selected odd prime dividing such an even
order would force that order to be at least \(2p>3u/2\).
Every selected prime is therefore coprime to every such group order,
and hence to every prescribed annihilating modulus. This proves
item three, including the fixed-modulus avoidance imposed before
prime selection.

The retained conjugate-order theorem supplies compatible conditions
putting the reference point in the primary formal group and congruent
to \(P\) at the opposite place for every selected prime.
The usual extra divisibility by all selected primes can be strengthened
at the unused prime to
\(h_0\equiv p_*\pmod{p_*^2}\). The preceding coprimality makes
this compatible with all earlier congruences. The fixed initial
modulus kills the component group of the compact cyclic real closure.
Every compatible progression refinement has the same real closure in
that component, so a positive \(h_0\) in the required real neighborhood
exists. It gives item four.
\end{proof}

\section{The unbounded-pole family}

\begin{theorem}\label{thm:unbounded}
Retain the arithmetic premises of \cite{Scalar}, including its group,
prime-pattern and intrinsic pole-realization inputs. For every integer
input and every fixed odd scalar parameter \(L\ge3\), there are
fixed roots, a fixed common slope, fixed contact data and fixed auxiliary
points \(U_i\), together with infinitely many complete rational source
tuples obtained by varying only \(H\) among the point coordinates.
There is a fixed split prime \(p_*\) such that, for arbitrarily large
integers \(r\), the two rational components of
\[
              L_S=\ell_8(U_S)/\ell_8(H)
\]
both have valuation \(1-r\) at \(p_*\). In particular their
rational logarithmic heights are at least \((r-1)\log p_*\),
although the input, root indices, root heights, contact arguments,
contact radicals and physical weights remain fixed.
\end{theorem}
\begin{proof}
Choose the complete prepared tuple supplied by Lemma~\ref{lem:dummy},
initially using \(H_0=[h_0]P\) and \(U_i=[n_i h_0]P\).
Freeze every root, the slope, and every \(U_i\) at these values.
Only the reference point \(H\) will change.

Let \(\mathcal T\) consist of all primes at which one of the two
fixed physical weights has odd valuation, together with two and the
source's finite preparation primes. The unused prime is not in
\(\mathcal T\). At the initial tuple the two \(b_j\) have positive
valuation at every odd-weight prime and satisfy the prepared local
conditions at two. Each of these is an open condition on \(H\),
because all the rational functions involved are defined and nonzero
at the initial point. Choose neighborhoods of \(H_0\) at the finitely
many places over \(\mathcal T\) preserving these conditions.
At the real places choose a neighborhood preserving the selected
inequalities \(0<b_j<1\), all scalar denominator inequalities,
and all point-algebra guard nonvanishing.

These finite neighborhoods are met by a fixed congruence
\(h\equiv h_0\pmod M\). We may take \(M\) prime to \(p_*\).
Indeed, the old preparation congruences have moduli supported on the
initial fixed modulus and powers of the old preparation primes.
The additional neighborhoods at actual selected contacts use their
reduction orders and powers of those contacts. The fixed avoidance
and Lemma~\ref{lem:dummy}(3) make every such factor prime to
\(p_*\). Include in \(M\) the order of the component group of
the real cyclic closure.

Write \(d_*\) and \(d_*'\) for the retained primary and opposite
annihilating moduli at the unused prime. Preserve the compatible
conditions on \(h\) supplied by the initial tuple, in particular
\([h]P\equiv O\) at the primary place and
\([h]P\equiv P\) at the opposite place. They are compatible with
\(h\equiv h_0\pmod M\). Their moduli are prime to \(p_*\).
For every \(r\ge2\), CRT therefore permits the further condition
\[
                         h\equiv p_*^r\pmod{p_*^{r+1}}.
\]
The same compact-group progression argument supplies a positive
integer \(h_r\) satisfying these conditions and the chosen real
neighborhood. Put \(H_r=[h_r]P\).

Let \(\mathfrak p_*\) be the primary place and let \(d\) be the
exact order of the reduction of \(P\) there. It is prime to \(p_*\).
Set \(e=v_{\mathfrak p_*}(z([d]P))\ge1\), where \(z=-x/y\).
The formal-group logarithm is an isometry at this good prime, and
\(\ell_8/z\) is a unit throughout its formal group. Since
\(v_{p_*}(n_S)=0\), the frozen point and varying reference give
\[
 v_{\mathfrak p_*}(\ell_8(U_S))=e+1,
 \qquad v_{\mathfrak p_*}(\ell_8(H_r))=e+r.
\]
Consequently \(v_{\mathfrak p_*}(L_S)=1-r\).
At the opposite place, \(H_r\) reduces to \(P\) and the frozen
\(U_S\) reduces to \(S\); both \(\ell_8\)-values are units by
the fixed avoidance, so \(L_S\) is a unit there.
If \(L_S=a_r+b_r\sqrt2\), half the sum and the scaled
half-difference of its two local values show
\[
                   v_{p_*}(a_r)=v_{p_*}(b_r)=1-r,
\]
because \(2\sqrt2\) is a unit and the two input valuations differ.
The height lower bounds follow from their reduced denominators.

It remains to verify that these tuples have complete pole witnesses.
At \(\mathcal T\) the neighborhoods preserve every odd-weight
valuation and the required two-adic class and depth. At every other
prime both weight valuations are even, so Lemma~\ref{lem:even}
preserves the square reduced denominators regardless of the new
auxiliary poles. The real inequalities remain in the prepared range.
Thus the intrinsic theorem supplies rational witness pairs for each
of the two actual values of \(b_j\).

At every real point-algebra character the roots and detector data
remain in the common guarded neighborhood. In particular the
detector factors are positive, their reciprocal traces are nonzero,
and the character values of \(b_j\) are nonzero. The intrinsic
witness density avoids all finitely many radicand and denominator
norms, exactly as in \cite{Scalar}. The existing finite-algebra
encoder supplies its rational coordinate. These are complete source
solutions, not merely local tuples.
\end{proof}

\begin{corollary}
No finite upper bound for the heights of the two rational components
of \(L_S\), depending only on the fixed input, root indices, root
heights, common slope and contact data, follows from the existing
source equations and guards. In particular, a proposed elimination
of a Pell point using a monic equation with coefficients depending on
\(L_S\) cannot justify such a coefficient-height bound by those
equations alone.
\end{corollary}

The theorem does not exclude an estimate on a smaller locus defined
by an additional monic Pell-cover equation. It does not apply to a
native-square source unless that additional rational-square condition
is also proved along the family. It gives no lower bound for the
number of variables in a definition of the integers.

\begin{thebibliography}{9}\raggedright
\bibitem{Scalar}\emph{Intrinsic congruent-curve poles from a fixed scalar
rescaling},
\nolinkurl{burungale_scalar_transfer_20261117_v56.tex}.
\bibitem{Precision}\emph{Scalar precision two with an eighth-order
elliptic comparison},
\nolinkurl{precision_transfer_20261117_v56.tex}.
\bibitem{Intrinsic}\emph{A two-coordinate pole test of coefficient height
four},
\nolinkurl{../single_cover_realization_v56/intrinsic_burungale_pole.tex}.
\bibitem{General}\emph{A defining polynomial in twenty-five variables},
\nolinkurl{papers/few_variables25/few_variables25_merged_20261110_v19.tex},
Section 6, especially the prime-cluster and progression-refinement
arguments.
\bibitem{Sun}Z.-W.~Sun, \emph{86 unknowns for Hilbert's tenth problem
over the rationals}, arXiv:2610.20691v2, Theorem 4.1.
\end{thebibliography}
\end{document}
